% ============================================================
% f2014_anhui_arts_q19_pyramid.tex
% 2014 安徽卷（文科）第 19 题：四棱锥 P-ABCD
% 来源 PDF：普通高考/2014/2014安徽文.pdf
% 原页：第 2 页（A3 横排）左栏下部，印面 860，第 19 题题干下方
% 说明：底面为边长 8 的正方形 ABCD，四条侧棱长 2√17，
%       顶点 P 在底面中心正上方，高为 6。
%       G，E，F，H 分别在 PB，AB，CD，PC 上，且 E 满足 EB=2，
%       平面 GEFH 垂直于底面，故其为过 EF（EF∥BC）的竖直平面。
%       被底面遮住的三条棱 PD、AD、DC 用虚线。
%       截面四边形 GEFH 中 GE、GH、HF 可见，EF 位于底面内被遮挡。
% ============================================================
\begin{tikzpicture}[>=Stealth, line width=0.7pt, font=\small,
  x=1cm, y=1cm, line cap=round, line join=round]

  % 底面斜二测画法：A(0,0), B(5.2,0), C(7.4,2.6), D(2.2,2.6)
  \coordinate (A) at (0,0);
  \coordinate (B) at (5.2,0);
  \coordinate (C) at (7.4,2.6);
  \coordinate (D) at (2.2,2.6);
  \coordinate (P) at (3.7,6.2);
  % E：AB 上 EB=2，AE=6
  \coordinate (E) at (3.9,0);
  % F：CD 上 CF=2
  \coordinate (F) at (6.1,2.6);
  % G，H 分别为 PB，PC 中点（平面为竖直平面 EF）
  \coordinate (G) at (4.45,3.1);
  \coordinate (H) at (5.55,4.4);

  % 不可见棱
  \draw[dashed] (P) -- (D);
  \draw[dashed] (A) -- (D);
  \draw[dashed] (D) -- (C);
  \draw[dashed] (E) -- (F);

  % 可见棱
  \draw (P) -- (A);
  \draw (P) -- (B);
  \draw (P) -- (C);
  \draw (A) -- (B);
  \draw (B) -- (C);

  % 截面四边形 GEFH（EF 已用虚线画出）
  \draw (G) -- (E);
  \draw (G) -- (H);
  \draw (H) -- (F);

  % 顶点标注
  \node[above] at (P) {$P$};
  \node[below left] at (A) {$A$};
  \node[below] at (B) {$B$};
  \node[right] at (C) {$C$};
  \node[above left, inner sep=2pt] at (D) {$D$};
  \node[below] at (E) {$E$};
  \node at (6.18,2.18) {$F$};
  \node at (4.02,3.10) {$G$};
  \node at (5.86,4.62) {$H$};

\end{tikzpicture}
